%---------------------------------------------------------------------------------------

\chapter[Considerações Finais]{Considerações Finais}\label{ch:consideracoes-finais}

%---------------------------------------------------------------------------------------

% Lista de abreviaturas e de siglas
\criarsigla{HIP}{\textit{Heterogeneous Interface for Portability}}

%---------------------------------------------------------------------------------------

% Lista de simbolos

%---------------------------------------------------------------------------------------

\section{Síntese dos Resultados Obtidos}\label{sec:sintese-dos-resultados-obtidos}

Este trabalho apresentou um miniaplicativo que implementa o método de Lattice Boltzmann para escoamentos compressíveis, validado no problema do tubo de choque de Sod e avaliado quanto à escalabilidade em CPUs multinúcleo e em aceleradores gráficos.
Os achados do Capítulo~\ref{ch:resultados} permitem responder de forma direta aos objetivos específicos enunciados na Seção~\ref{subsec:objetivos-especificos}.

Quanto ao primeiro objetivo específico, relativo à validação da implementação, os perfis de massa específica, de velocidade e de pressão obtidos acompanham a solução analítica do problema de Riemann nas regiões suaves do escoamento e posicionam corretamente a cabeça e a cauda da rarefação, a descontinuidade de contato e a frente de choque.
O estudo de refinamento indicou taxa de convergência efetiva próxima da primeira ordem em normas globais, resultado esperado e coerente com a presença de descontinuidades que dominam a norma do erro, ainda que o esquema de propagação seja formalmente de segunda ordem em regiões contínuas.
Os estudos de sensibilidade acrescentaram a esse quadro dois limites práticos de operação: o crescimento monotônico do erro com a razão de pressões inicial, até a perda de estabilidade nas proximidades de $\Pi = 20$, e a existência de uma faixa intermediária de viscosidade que minimiza o erro, situada entre o excesso de dissipação e a instabilidade associada ao limite inferior do tempo de relaxação.

Quanto ao segundo objetivo específico, referente à análise de desempenho das diferentes filosofias de programação, verificou-se que as tecnologias voltadas a CPU multinúcleo convergem para acelerações da mesma ordem de grandeza no regime de malhas grandes, o que indica que a escolha entre elas afeta primariamente a portabilidade e o esforço de desenvolvimento, e apenas secundariamente o desempenho alcançável.
O uso de aceleradores gráficos, por sua vez, proporciona ganhos de ordem de grandeza superior, condicionados a malhas suficientemente refinadas para amortizar a latência de lançamento dos \textit{kernels} e a transferência inicial de dados, o que evidencia o potencial das arquiteturas heterogêneas para o LBM compressível.

Quanto ao terceiro objetivo específico, relativo à escalabilidade, a investigação demonstrou que o algoritmo é limitado primordialmente pela largura de banda de memória, e não pela capacidade de cálculo aritmético.
A eficiência paralela permanece elevada enquanto o conjunto de trabalho se acomoda nos níveis mais rápidos da hierarquia de memória e degrada-se à medida que a malha cresce, quando a contenção do barramento compartilhado entre os núcleos passa a dominar o comportamento observado, efeito que a lei de Amdahl, isoladamente, não explica.

Quanto ao quarto objetivo específico, que confronta desempenho e custo de desenvolvimento, a contagem de linhas de código mostrou que as tecnologias baseadas em diretivas oferecem a melhor relação de compromisso quando o prazo de desenvolvimento é restritivo, ao passo que os modelos de baixo nível, como o CUDA e o OpenCL, exigem investimento consideravelmente maior em código e em manutenção para extrair o desempenho de pico.
As camadas de abstração situam-se em posição intermediária e mostram-se a escolha mais equilibrada quando o mesmo código precisa atender simultaneamente a CPUs e a aceleradores de fabricantes distintos.

%---------------------------------------------------------------------------------------

\section{Contribuições do Trabalho}\label{sec:contribuicoes-do-trabalho}

Em termos de contribuição acadêmica, o estudo fornece uma linha de base reprodutível para comparações entre operadores de colisão, discretizações do espaço de velocidades e estratégias de filtragem ou de dissipação controlada no LBM compressível.
Apresenta, além disso, um perfil de desempenho que explicita gargalos práticos, como a largura de banda de memória e as sobrecargas de sincronização, o que orienta pesquisas futuras em otimização de \textit{kernels} em arquiteturas heterogêneas.
Destaca-se também o desenvolvimento de um miniaplicativo estruturado, que serve como artefato experimental útil para o ensino de computação de alto desempenho e para a prototipagem de novos modelos físicos.

%---------------------------------------------------------------------------------------

\section{Limitações Identificadas}\label{sec:limitacoes-identificadas}

Reconhece-se honestamente que este trabalho apresenta limitações decorrentes do ambiente experimental e do estágio atual de desenvolvimento do resolvedor.
Os ensaios foram conduzidos em uma única máquina, equipada com um processador Intel Core 7 240H e uma GPU dedicada, sem acesso a agrupamentos de computadores (\textit{clusters}) que permitiriam avaliar a escalabilidade distribuída em larga escala via MPI.
Adicionalmente, parte das comparações entre tecnologias em GPU para malhas extremamente refinadas permanece como expectativa teórica a confirmar, dada a necessidade de recursos de memória de vídeo superiores aos disponíveis no sistema de ensaio.
A ausência de testes com descontinuidades mais severas ou escoamentos multidimensionais também delimita o escopo da validação realizada.

%---------------------------------------------------------------------------------------

\section{Trabalhos Futuros}\label{sec:trabalhos-futuros}

Como trabalhos futuros, propõe-se um roteiro de evolução estruturado em quatro frentes, cada uma delas respondendo a uma das limitações reconhecidas na seção anterior.

Na frente de portabilidade de desempenho, sugere-se estender as implementações em Kokkos e em RAJA a espaços de execução ainda não avaliados, em particular HIP e SYCL, de modo a verificar se as conclusões obtidas para GPUs NVIDIA se mantêm em aceleradores de outros fabricantes.
Recomenda-se, no mesmo esforço, investigar disposições de dados alternativas, com blocagem espacial, alinhamento explícito e busca antecipada, que possam reduzir a pressão sobre o subsistema de memória identificada como fator limitante do desempenho.

Na frente de paralelismo distribuído, propõe-se estender a decomposição de domínio unidimensional atualmente implementada para uma decomposição multidimensional, aliada à agregação de mensagens para mitigação de latência e a estratégias de balanceamento de carga para partições desiguais em sistemas heterogêneos.
Propõe-se ainda avaliar a combinação da MPI com as implementações voltadas a aceleradores, configuração híbrida que explora simultaneamente o paralelismo de grão fino no interior de cada dispositivo e o paralelismo de grão grosso entre nós de computação.

Na frente experimental, faz-se necessária uma campanha extensiva de ensaios em ambiente de computação de alto desempenho, incluindo estudos de escalabilidade forte e fraca em agrupamentos equipados com aceleradores e com nós de CPU de maior contagem de núcleos, variando o refinamento espacial e o número de passos de tempo.
Nessas campanhas, devem ser medidos a taxa de atualização de nós, a eficiência paralela, o posicionamento em relação ao modelo \textit{roofline} e o custo energético, métrica cuja relevância cresce nos grandes centros de computação.

Na frente física e numérica, recomenda-se a validação com descontinuidades mais severas, entre elas os problemas de Lax e de Shu-Osher, confrontadas com soluções exatas ou de referência, e a extensão dos ensaios a configurações bidimensionais e tridimensionais, nas quais o padrão de acesso à memória da etapa de propagação torna-se substancialmente mais exigente.
Estudos de sensibilidade a parâmetros numéricos, como a intensidade da regularização e o número de iterações do resolvedor local de temperatura, permitirão ainda mapear o compromisso entre estabilidade e custo computacional.

Este trabalho conclui que o miniaplicativo alcançou um grau de maturidade que justifica sua aplicação em infraestruturas de computação de alto desempenho de maior porte.
A trajetória percorrida posiciona o LBM compressível como uma ferramenta promissora para investigações que demandam alta fidelidade física e eficiência computacional em arquiteturas heterogêneas.
